\section{Degree-corrected models and oracle benchmarks}\label{sec:statements}
The target is a constant oracle-risk bound
\begin{equation}\label{eq:target}
 \E\ell(\widehat z,z)\le C R_*(n),
\end{equation}
with $C$ independent of $n$. The companion retains ordinary-SBM
notation where needed by the historical free-profile analysis.
The independent ordinary-SBM article contains the current ordinary
adjacency block-LRT theorem and its complete proof.

\subsection{Bernoulli models and fixed structural constants}

There are $n$ vertices and $K\ge2$ groups, with fixed $K$.
Let $z\in[K]^n$, let $Z$ be its indicator matrix, and let $n_a$ be the size
of group $a$. Throughout $n_a\ge\pi_-n$ for a fixed $\pi_->0$.
The edges above the diagonal are independent, the adjacency matrix $A$ is
symmetric, and $A_{ii}=0$. In DC-SBM,
\begin{equation}\label{eq:mean-model}
 A_{ij}\sim\Ber(\theta_i\theta_j p C_{z_i z_j}),\quad i<j,
 \qquad d=np,\qquad
 \Omega=\Theta Z(pC)Z^\top\Theta.
\end{equation}
Here $\Theta=\diag(\theta_i)$, $C=C^\top$, and the fixed structural bounds are
\begin{equation}\label{eq:structural}
 0<\theta_-\le\theta_i\le\theta_+<\infty,\quad
 \sum_{i:z_i=a}\theta_i=n_a,\quad
 0<c_-\le C_{ab}\le c_+<\infty,\quad
 \sigma_{\min}(C)\ge\kappa>0.
\end{equation}
The model class is nonempty only when $1\in[\theta_-,\theta_+]$.
All constants in asymptotic bounds may depend on these structural bounds,
but not on $n$. One may take a fixed $C$, or a sequence in this same compact
class. Ordinary SBM is the special case $\theta_i=1$, with block probability
matrix $P=pC$.

The completed matrix $\Omega$ includes a diagonal for algebraic convenience;
the actual mean is $\E A=\Omega-\diag(\Omega_{ii})$.
That deterministic difference has norm $O(p)$ and is included in all error
bounds. Put $D=\diag(A\one)$ and
$\lambda_i=(\E A\one)_i$. Thus $\lambda_i\asymp d$ and
$\Omega\one$ equals $d\theta_i S_{z_i}$ in coordinate $i$, where
\begin{equation}\label{eq:exposures}
 \pi_a=n_a/n,\qquad S_a=\sum_b\pi_b C_{ab}.
\end{equation}
The loss for the estimated partition is
\[
 \ell(\widehat z,z)=\frac1n\min_{\sigma\in\mathfrak S_K}
 \sum_i\ind{\sigma(\widehat z_i)\ne z_i}.
\]

\subsection{Three explicitly different local oracles}\label{sec:oracle}

All oracle classifiers below use the raw incident row $A_{i,-i}$ and the
true neighbor labels. They are local likelihood classifiers: an exact global
block-size constraint is not used to determine the one missing label.
For simplicity candidate weights are uniform. Any fixed weights bounded
above and below by positive constants give the same conclusions.

For ordinary SBM, the oracle chooses the maximizer of
\begin{equation}\label{eq:sbm-oracle}
 \ell_{i,*}(a)=\sum_{j\ne i}
 \{A_{ij}\log P_{a z_j}+(1-A_{ij})\log(1-P_{a z_j})\}.
\end{equation}
Its average error probability is $R_{\mathrm{SBM},*}$.
For DC-SBM, the \emph{known-target-degree oracle} replaces $P_{a z_j}$ in
\eqref{eq:sbm-oracle} by $\theta_i\theta_j p C_{a z_j}$ and is supplied
with every true $\theta_j$, including $\theta_i$. Its risk is
$R_{\theta,*}$.

The \emph{profile oracle} is supplied with true neighbor parameters and
$pC$, but profiles the target degree separately for every candidate:
\begin{equation}\label{eq:profile-oracle}
 \ell_{i,\mathrm{prof},*}(a)=
 \sup_{t\ge0}\sum_{j\ne i}
 \{A_{ij}\log(t\theta_jpC_{a z_j})
 +(1-A_{ij})\log(1-t\theta_jpC_{a z_j})\}.
\end{equation}
We restrict candidate edge probabilities to at most $q_{\max}=1/2$; the
constraint is inactive on the events used in the sparse theorem. Write
$R_{\mathrm{prof},*}$ for its average error probability. This is a
full-edge generalized likelihood-ratio reference. It is not described as a
Bayes test between two fixed known-$\theta_i$ laws. In all three definitions,
the risk averages the actual directional error at each fixed true vertex.
Ties are resolved by a fixed ordering; strict-tail arguments below cover
any such ordering.

